% Minimal TeXtured Example
% This demonstrates the basic usage of the TeXtured document class

% IMPORTANT: PDF/A compliance requires \DocumentMetadata before \documentclass
\DocumentMetadata{lang=en, pdfversion=1.7, pdfstandard=A-2u}

% Pass options to the class using \PassOptionsToClass
\PassOptionsToClass{
    FANCY,       % Enable fancy typography
    % WIP,       % Work-in-progress mode (todos, draft watermarks)
    noLINKBOXES, % Disable colored link boxes (faster compilation)
    % CENSOR,    % Enable censoring of designated passages
}{textured}

\documentclass[12pt,a4paper]{textured}

% Configure document metadata using key-value interface
\SetKeys[textured]{
  thesistype       = {master},           % bachelor, master, or doctoral
  title            = {A Minimal TeXtured Example},
  titleplaintext   = {A Minimal TeXtured Example},
  author           = {Your Name},
  authorplaintext  = {Your Name},
  yearsubmitted    = {2025},
  department       = {Department of Computer Science},
  depttype         = {Department},       % Department or Institute
  supervisor       = {Prof. Advisor Name},
  supervisorsdepartment = {Supervisor's Department},
  studyprogramme   = {Computer Science},
  keywords         = {LaTeX, thesis, TeXtured, example},
  keywordsplaintext = {LaTeX, thesis, TeXtured, example},
  subject          = {Demonstration of TeXtured class features},
  abstract         = {This is a minimal example demonstrating the TeXtured document class for academic theses.},
}

% Bibliography: Add your own bibliography files
% \addbibresource{mybibliography.bib}

% Optional: To cite TeXtured itself, the package provides textured.bib
% (After CTAN installation, this will be findable automatically)
\addbibresource{textured.bib}

\begin{document}

\frontmatter

% Optional: Use included MFF CUNI frontmatter templates
% \include{frontmatter/title}
% \include{frontmatter/declaration}
% \include{frontmatter/dedication}

% Or create custom title page:
\begin{titlepage}
\centering
\vspace*{2cm}
{\Huge\bfseries A Minimal TeXtured Example\par}
\vspace{1.5cm}
{\Large Your Name\par}
\vspace{0.5cm}
{\large Department of Computer Science\par}
\vspace{3cm}
{\large A demonstration of the TeXtured document class\par}
\vfill
{\large Master's Thesis\par}
\vspace{0.5cm}
{\large January 2025\par}
\end{titlepage}

\cleardoublepage

\cleardoublepage

\tableofcontents

\cleardoublepage

\mainmatter

\chapter{Introduction}

This is a minimal example demonstrating the TeXtured document class, a professionally typeset \LaTeX{} class for academic theses.

\section{What is TeXtured?}

TeXtured provides a complete framework for writing elegant academic documents with:
\begin{itemize}
  \item Professional typography with optional fancy features
  \item Modern \LaTeX{} conventions (\verb|\NewDocumentCommand|, \verb|\DeclareKeys|)
  \item PDF/A-2u compliance for long-term archival
  \item Rich mathematical typesetting and theorem environments
  \item Work-in-progress mode with todos and debug features
  \item Modular architecture with specialized style packages
\end{itemize}

\section{Quick Example}

To use TeXtured, simply load the class and configure your document metadata:

\begin{verbatim}
\DocumentMetadata{pdfstandard=A-2u}
\PassOptionsToClass{FANCY}{textured}
\documentclass{textured}
\SetKeys[textured]{
  title  = {My Thesis},
  author = {Your Name},
  ...
}
\end{verbatim}

For complete documentation, see the TeXtured manual and the repository at \url{https://github.com/jdujava/TeXtured}.

\chapter{Mathematical Content}

TeXtured provides excellent support for mathematical typesetting.

\section{Equations}

Here is a simple equation:
\begin{equation}
  E = mc^2
  \label{eq:einstein}
\end{equation}

And a more complex example using aligned equations:
\begin{align}
  \nabla \times \mathbf{E} &= -\frac{\partial \mathbf{B}}{\partial t}, \\
  \nabla \times \mathbf{B} &= \mu_0 \mathbf{J} + \mu_0 \varepsilon_0 \frac{\partial \mathbf{E}}{\partial t}.
\end{align}

\section{Theorem Environments}

TeXtured includes pre-configured theorem environments:

\begin{theorem}[Pythagorean Theorem]
  \label{thm:pythagoras}
  For a right triangle with legs $a$ and $b$, and hypotenuse $c$:
  \begin{equation}
    a^2 + b^2 = c^2 \eqend
  \end{equation}
\end{theorem}

\begin{proof}
  The proof is left as an exercise for the reader.
\end{proof}

\begin{lemma}
  Every result used in a proof should be stated as a lemma.
\end{lemma}

\begin{corollary}
  Theorem~\ref{thm:pythagoras} implies that the hypotenuse is always the longest side.
\end{corollary}

\chapter{Advanced Features}

\section{Remark Environments}

\begin{definition}
  A \emph{minimal example} is the smallest complete demonstration of a concept.
\end{definition}

\begin{remark}
  TeXtured provides many remark-like environments for supplementary content.
\end{remark}

\begin{example}
  Consider the triangle with sides $a=3$, $b=4$, and $c=5$. By Theorem~\ref{thm:pythagoras}:
  \[
    3^2 + 4^2 = 9 + 16 = 25 = 5^2
  \]
  Hence this is a right triangle.
\end{example}

\section{Citing TeXtured}

If you use TeXtured for your thesis, you can cite it using the provided bibliography entry \autocite{TeXtured}.

\section{Further Reading}

For more examples and complete documentation:
\begin{itemize}
  \item See \texttt{thesis.tex} in the repository for a full-featured example
  \item Read the TeXtured manual (PDF documentation)
  \item Visit the GitHub repository: \url{https://github.com/jdujava/TeXtured}
\end{itemize}

\backmatter

\references

\end{document}
